\documentclass[a4paper, 12pt]{extarticle}


%%%%%%%%%%%%%%%%%%%%%%%%%%%%%%
% PACKAGES
%%%%%%%%%%%%%%%%%%%%%%%%%%%%%%

%!TEX encoding = UTF8

%%%%%%%%%%%%%%%%%%%%%%%%%%%%%%%%%
% PACKAGE IMPORTS
%%%%%%%%%%%%%%%%%%%%%%%%%%%%%%%%%


\usepackage[french]{babel}
\usepackage[T1]{fontenc}
\usepackage[utf8]{inputenc}

\usepackage[
	margin=2cm,
	right=2.2cm,
	footskip=1.5cm,
]{geometry}

%ams
\usepackage{amsmath,amsfonts,amsthm,amssymb,mathtools}

% for fancy headers
\usepackage{fancyhdr}

% for table of content stuff
\usepackage{titlesec}

\usepackage{bookmark}
\usepackage{enumitem}
\usepackage[most,many,breakable]{tcolorbox}
\usepackage{varwidth,etoolbox}
\usepackage[makeroom]{cancel}
\usepackage{xcolor}
\usepackage{multicol,array}
% algorithms
\usepackage[ruled,vlined,linesnumbered]{algorithm2e}

% delays answers
% answers are from last exercise if no label is provided
% delays exercices because of exercise bank (exercices.tex file in course folder)
%\usepackage[answerdelayed, exercisedelayed, lastexercise]{exercise}
\usepackage[answerdelayed, lastexercise]{exercise}

% for slanted sfrac
\usepackage{xfrac}
% for \uline which allows underline with line breaks
\usepackage[normalem]{ulem}

\usepackage{caption, subcaption}

% margins for points (and old (p. %page%) ).
\usepackage{marginnote}

\usepackage{makecell} % pour \thead 

%index
\usepackage[texindy]{imakeidx}
\makeindex[
	columns=2, 
	title=Index, 
	options= -L french,
	%intoc,
]
\indexsetup{level=\chapter}

%virgules
\usepackage{icomma}

% for striked out implies sign (\centernot\implies)
\usepackage{centernot}

% roman numerals for \section
\renewcommand{\thesection}{\Roman{section}} 

% tikz
\usepackage{tikz}

% package systeme pour les systèmes d'équations bien alignés
\usepackage{systeme}
% add \alpha \beta \gamma as variables
% put \\ as separator so commas don't return to line
\setsysteme{eq sep = \\, align={r,r}}
% example : 
% \systeme[sort={\alpha,\beta, \gamma}]{ \alpha + 2\beta = 4, \\ \beta = 8.} 

% for polynomial long division
\usepackage{polynom}
% example : 
% \polylongdiv[style=D]{x^3 - 1 + 2x - 2x^2}{x-1}


% date
\usepackage{advdate}


%--------------------------------------------------------------------
\makeatletter
%                            \fillwithlines


% \fillwithlines takes one argument, which is either a length or \fill
% or \stretch{number}, and it fills that much vertical space with
% horizontal lines that run the length of the current line.  That is,
% they extend from the current left margin (which depends on whether
% we're in a question, part, subpart, or subsubpart) to the right
% margin.
%
% The distance between the lines is \linefillheight, whose default value
% is set with the command
%
% \setlength\linefillheight{.25in}
%
% This value can be changed by giving a new \setlength command.
%
% The thickness of the lines is \linefillthickness, whose default value
% is set with the command
%
% \setlength\linefillthickness{.1pt}
%
% This value can be changed by giving a new \setlength command.


\newlength\linefillheight
\newlength\linefillthickness
\setlength\linefillheight{.25in}
\setlength\linefillthickness{0.1pt}

\newcommand\linefill{\leavevmode
    \leaders\hrule height \linefillthickness \hfill\kern\z@}


\def\fillwithlines#1{%
  \begingroup
  \ifhmode
    \par
  \fi
  \hrule height \z@
  \nobreak
  \setbox0=\hbox to \hsize{\hskip \@totalleftmargin
          \vrule height \linefillheight depth \z@ width \z@
          \linefill}%
  % We use \cleaders (rather than \leaders) so that a given
  % vertical space will always produce the same number of lines
  % no matter where on the page it happens to start:
  \cleaders \copy0 \vskip #1 \hbox{}%
  \endgroup
}
\makeatother
%--------------------------------------------------------------------


% tikz
\usepackage{tikz, pgfplots}
\usepackage{prerex}
% i wish external worked but idk it sucks
%\usetikzlibrary{external}
%\tikzexternalize[prefix=figures/]

% for function graph
\usetikzlibrary{positioning}
\usetikzlibrary{shapes.geometric}
\usetikzlibrary{positioning}
\usetikzlibrary{shapes.misc}
\tikzset{
dot/.style = {circle, fill=#1, minimum size=5pt,
              inner sep=0pt, outer sep=0pt},
dot/.default = black % size of the circle diameter
}
\tikzset{cross/.style={cross out, draw=black, minimum size=2*(#1-\pgflinewidth), inner sep=0pt, outer sep=0pt},
%default radius will be 1pt. 
cross/.default={1pt}}

 % for braces
\usetikzlibrary{decorations.pathreplacing}
% for hashing area
\usetikzlibrary{patterns}
% tableaux var, signe
% source https://www.sqlpac.com/fr/documents/latex-package-tkz-tab-tikz-tableaux-de-signes-et-de-variations-de-fonctions.html
\usepackage{tkz-tab}



\tikzset{
	every axis/.style = {clip=true, grid style = {opacity=.5}}
}


\pagestyle{fancy}
\fancyhead[L]{Université de Mayotte}
\fancyhead[C]{\bf Contrôle Continu}
\fancyhead[R]{\today}

%%%%%%%%%%%%%%%%%%%%%%%%%%%%%%
% SELF MADE COLORS
%%%%%%%%%%%%%%%%%%%%%%%%%%%%%%

%!TEX encoding = UTF8
%%%%%%%%%%%%%%%%%%%%%%%%%%%%%%
% SELF MADE COLORS
%%%%%%%%%%%%%%%%%%%%%%%%%%%%%%


\definecolor{myg}{RGB}{56, 140, 70}
\definecolor{myb}{RGB}{45, 111, 177}
\definecolor{myr}{RGB}{199, 68, 64}
\definecolor{mytheorembg}{HTML}{F2F2F9}
\definecolor{mytheoremfr}{HTML}{00007B}
\definecolor{mylenmabg}{HTML}{FFFAF8}
\definecolor{mylenmafr}{HTML}{983b0f}
\definecolor{mypropbg}{HTML}{f2fbfc}
\definecolor{mypropfr}{HTML}{191971}
\definecolor{myexamplebg}{HTML}{F2FBF8}
\definecolor{myexamplefr}{HTML}{88D6D1}
\definecolor{myexampleti}{HTML}{2A7F7F}
\definecolor{mydefinitbg}{HTML}{E5E5FF}
\definecolor{mydefinitfr}{HTML}{3F3FA3}
\definecolor{notesgreen}{RGB}{0,162,0}
\definecolor{myp}{RGB}{197, 92, 212}
\definecolor{mygr}{HTML}{2C3338}
\definecolor{myred}{RGB}{127,0,0}
\definecolor{myyellow}{RGB}{169,121,69}
\definecolor{myexercisebg}{HTML}{F2FBF8}
\definecolor{myexercisefg}{HTML}{88D6D1}
\definecolor{doc}{RGB}{0,60,110}

% manim colors because they're beautiful
% https://docs.manim.community/en/stable/reference/manim.utils.color.manim_colors.html

\definecolor{BLACK}{HTML}{000000}\definecolor{BLUE}{HTML}{58C4DD}\definecolor{BLUE_A}{HTML}{C7E9F1}\definecolor{BLUE_B}{HTML}{9CDCEB}\definecolor{BLUE_C}{HTML}{58C4DD}\definecolor{BLUE_D}{HTML}{29ABCA}\definecolor{BLUE_E}{HTML}{236B8E}\definecolor{DARKER_GRAY}{HTML}{222222}\definecolor{DARKER_GREY}{HTML}{222222}\definecolor{DARK_BLUE}{HTML}{236B8E}\definecolor{DARK_BROWN}{HTML}{8B4513}\definecolor{DARK_GRAY}{HTML}{444444}\definecolor{DARK_GREY}{HTML}{444444}\definecolor{GOLD}{HTML}{F0AC5F}\definecolor{GOLD_A}{HTML}{F7C797}\definecolor{GOLD_B}{HTML}{F9B775}\definecolor{GOLD_C}{HTML}{F0AC5F}\definecolor{GOLD_D}{HTML}{E1A158}\definecolor{GOLD_E}{HTML}{C78D46}\definecolor{GRAY}{HTML}{888888}\definecolor{GRAY_A}{HTML}{DDDDDD}\definecolor{GRAY_B}{HTML}{BBBBBB}\definecolor{GRAY_BROWN}{HTML}{736357}\definecolor{GRAY_C}{HTML}{888888}\definecolor{GRAY_D}{HTML}{444444}\definecolor{GRAY_E}{HTML}{222222}\definecolor{GREEN}{HTML}{83C167}\definecolor{GREEN_A}{HTML}{C9E2AE}\definecolor{GREEN_B}{HTML}{A6CF8C}\definecolor{GREEN_C}{HTML}{83C167}\definecolor{GREEN_D}{HTML}{77B05D}\definecolor{GREEN_E}{HTML}{699C52}\definecolor{GREY}{HTML}{888888}\definecolor{GREY_A}{HTML}{DDDDDD}\definecolor{GREY_B}{HTML}{BBBBBB}\definecolor{GREY_BROWN}{HTML}{736357}\definecolor{GREY_C}{HTML}{888888}\definecolor{GREY_D}{HTML}{444444}\definecolor{GREY_E}{HTML}{222222}\definecolor{LIGHTER_GRAY}{HTML}{DDDDDD}\definecolor{LIGHTER_GREY}{HTML}{DDDDDD}\definecolor{LIGHT_BROWN}{HTML}{CD853F}\definecolor{LIGHT_GRAY}{HTML}{BBBBBB}\definecolor{LIGHT_GREY}{HTML}{BBBBBB}\definecolor{LIGHT_PINK}{HTML}{DC75CD}\definecolor{LOGO_BLACK}{HTML}{343434}\definecolor{LOGO_BLUE}{HTML}{525893}\definecolor{LOGO_GREEN}{HTML}{87C2A5}\definecolor{LOGO_RED}{HTML}{E07A5F}\definecolor{LOGO_WHITE}{HTML}{ECE7E2}\definecolor{MAROON}{HTML}{C55F73}\definecolor{MAROON_A}{HTML}{ECABC1}\definecolor{MAROON_B}{HTML}{EC92AB}\definecolor{MAROON_C}{HTML}{C55F73}\definecolor{MAROON_D}{HTML}{A24D61}\definecolor{MAROON_E}{HTML}{94424F}\definecolor{ORANGE}{HTML}{FF862F}\definecolor{PINK}{HTML}{D147BD}\definecolor{PURE_BLUE}{HTML}{0000FF}\definecolor{PURE_GREEN}{HTML}{00FF00}\definecolor{PURE_RED}{HTML}{FF0000}\definecolor{PURPLE}{HTML}{9A72AC}\definecolor{PURPLE_A}{HTML}{CAA3E8}\definecolor{PURPLE_B}{HTML}{B189C6}\definecolor{PURPLE_C}{HTML}{9A72AC}\definecolor{PURPLE_D}{HTML}{715582}\definecolor{PURPLE_E}{HTML}{644172}\definecolor{RED}{HTML}{FC6255}\definecolor{RED_A}{HTML}{F7A1A3}\definecolor{RED_B}{HTML}{FF8080}\definecolor{RED_C}{HTML}{FC6255}\definecolor{RED_D}{HTML}{E65A4C}\definecolor{RED_E}{HTML}{CF5044}\definecolor{TEAL}{HTML}{5CD0B3}\definecolor{TEAL_A}{HTML}{ACEAD7}\definecolor{TEAL_B}{HTML}{76DDC0}\definecolor{TEAL_C}{HTML}{5CD0B3}\definecolor{TEAL_D}{HTML}{55C1A7}\definecolor{TEAL_E}{HTML}{49A88F}\definecolor{WHITE}{HTML}{FFFFFF}\definecolor{YELLOW}{HTML}{FFFF00}\definecolor{YELLOW_A}{HTML}{FFF1B6}\definecolor{YELLOW_B}{HTML}{FFEA94}\definecolor{YELLOW_C}{HTML}{FFFF00}\definecolor{YELLOW_D}{HTML}{F4D345}\definecolor{YELLOW_E}{HTML}{E8C11C}

%%%%%%%%%%%%%%%%%%%%%%%%%%%%
% ENVIRONMENTS
%%%%%%%%%%%%%%%%%%%%%%%%%%%%

\theoremstyle{definition}

\newtheorem{theorem}{Théorème}
\newtheorem{corollaire}[theorem]{Corollaire}
\newtheorem{lemme}[theorem]{Lemme}
\newtheorem{proposition}[theorem]{Proposition}
\newtheorem{exercice}[theorem]{Exercice}
\newtheorem{exemple}[theorem]{Exemple}
\newtheorem{definition}[theorem]{Définition}
\newtheorem*{question}{Question}
\newtheorem*{preuve}{Preuve}
\newtheorem*{remarque}{Remarque}
\newtheorem*{strategie}{Stratégie}
\newtheorem*{methode}{Méthode}
\newtheorem*{notation}{Notation}
\newtheorem*{nomenclature}{Nomenclature}
\newtheorem{axiome}[theorem]{Axiome}
\newtheorem*{heuristique}{Heuristique}
\newtheorem*{motivation}{Motivation}

\newtheorem*{definition*}{Définition}
\newtheorem*{lemme*}{Lemme}
\newtheorem*{proposition*}{Proposition}
\newtheorem*{theorem*}{Théorème}
\newtheorem*{corollaire*}{Corollaire}

%%%%%%%%%%%%%%%%%%%%%%%%%%%%
% LETTERFONTS
%%%%%%%%%%%%%%%%%%%%%%%%%%%%

%!TEX encoding = UTF8

%fonts
\usepackage{libertinus,libertinust1math}
\usepackage[T1]{fontenc}

% for bold face math mode in section titles
\AddToHook{bfseries}{\boldmath}

% for calligraphic C, D, P (important to import this after the font)
\usepackage{calrsfs}
\newcommand{\D}{\mathcal{D}}
\newcommand{\Ccal}{\mathcal{C}}
\renewcommand{\P}{\mathcal{P}}

% Schwartz
\renewcommand{\S}{\mathcal{S}} % \S est le signe paragraphe normalement

% corps
\newcommand{\C}{\mathbb{C}}
\newcommand{\R}{\mathbb{R}}
\newcommand{\Rnn}{\mathbb{R}^{2n}}
\newcommand{\Z}{\mathbb{Z}}
\newcommand{\N}{\mathbb{N}}
\newcommand{\Q}{\mathbb{Q}}
\newcommand{\E}{\mathbb{E}}
\newcommand{\DD}{\mathbb{D}}
\newcommand{\K}{\mathbb{K}}
\newcommand{\F}{\mathbb{F}}
\newcommand{\U}{\mathbb{U}}
\renewcommand{\H}{\mathbb{H}}

% order notations
\DeclareRobustCommand{\O}{%
  \text{\usefont{OMS}{cmsy}{m}{n}O}%
}

% japanese bracket
\newcommand{\japb}[1]{\langle #1 \rangle}

% arrows over partial derivatives
\newcommand{\lp}{\overleftarrow{\partial}}
\newcommand{\rp}{\overrightarrow{\partial}}

% quantization
\newcommand{\h}{\hbar}
\newcommand{\Opht}{\textrm{Op}_{\h}^{t}}
\newcommand{\Op}[2][\hbar]{\textrm{Op}_{#1}^{#2}}

% omega functions
\newcommand{\omegap}[2][\rho_0]{\omega(\partial_{#1},\partial_{#2})}
\newcommand{\omegar}[2][\rho_0]{\omega(#1,#2)}

% space before semicolon
\mathcode`\;="303B

% for \Lightning
\usepackage{marvosym}

% for \warning
% 66 or 49 idk ; depends on the computer for some reason
\newcommand{\warning}{{\fontencoding{U}\fontfamily{futs}\selectfont\char 49\relax}}

% Q(\sqrt(d)) field
\newcommand{\Qsqrt}[1]{\Q\bigl(\mspace{-3mu}\sqrt{#1}\bigr)}

% closed off square root from https://en.wikibooks.org/wiki/LaTeX/Mathematics#Roots
% simple version: does not allow cbrt{x} = sqrt[3]{x}
% New definition of square root:
% it renames \sqrt as \oldsqrt
\let\oldsqrt\sqrt
% it defines the new \sqrt in terms of the old one
\def\sqrt{\mathpalette\DHLhksqrt}
\def\DHLhksqrt#1#2{%
\setbox0=\hbox{$#1\oldsqrt{#2\,}$}\dimen0=\ht0
\advance\dimen0-0.2\ht0
\setbox2=\hbox{\vrule height\ht0 depth -\dimen0}%
{\box0\lower0.4pt\box2}}



%%%%%%%%%%%%%%%%%%%%%%%%%%%%
% MACROS
%%%%%%%%%%%%%%%%%%%%%%%%%%%%

%!TEX encoding = UTF8


%%%%%%%%%%%%%%%%%%%%%%%%%%%%
% THEOREMS
%%%%%%%%%%%%%%%%%%%%%%%%%%%%

\newcommand{\thm}[3]{\begin{theorem}[#1]\label{#3}#2\end{theorem}}
\newcommand{\cor}[3]{\begin{corollaire}[#1]\label{#3}#2\end{corollaire}}
\newcommand{\lem}[3]{\begin{lemme}[#1]\label{#3}#2\end{lemme}}
\newcommand{\prop}[3]{\begin{proposition}[#1]\label{#3}#2\end{proposition}}
\newcommand{\ex}[3]{\begin{exemple}[#1]\label{#3}#2\end{exemple}}
\newcommand{\dfn}[3]{\begin{definition}[#1]\label{#3}#2\end{definition}}
\newcommand{\qs}[2]{\begin{question}[#1]#2\end{question}}

\renewcommand\qedsymbol{$\blacksquare$}
\newcommand{\pf}[2]{\begin{preuve}[#1]#2\hfill\qedsymbol\end{preuve}}

\newcommand{\nt}[1]{\begin{remarque}#1\end{remarque}}
\newcommand{\str}[1]{\begin{strategie}#1\end{strategie}}
\newcommand{\mth}[1]{\begin{methode}#1\end{methode}}
\newcommand{\ax}[3]{\begin{axiome}[#1]\label{#3}#2\end{axiome}}
\newcommand{\notations}[1]{\begin{notation}#1 \end{notation}}
\newcommand{\nomen}[1]{\begin{nomenclature}#1 \end{nomenclature}}
\newcommand{\heur}[1]{\begin{heuristique}#1\end{heuristique}}
\newcommand{\motiv}[1]{\begin{motivation}#1\end{motivation}}


%%%%%%%%%%%%%%%%%%%%%%%%%%%%
% EXERCISES
%%%%%%%%%%%%%%%%%%%%%%%%%%%%

\newcommand{\exe}[4]{
	\begin{Exercise}[#1, label=#3]
		%\marginpar{\mbox{\scriptsize(solution p.\pageref{\ExerciseLabel-Answer})}}
		#2
	\end{Exercise}
	\begin{Answer}[ref=#3]
		#4
	\end{Answer}
}

% version 1
\newcommand{\exemulticols}[5]{
	\begin{multicols}{2}
	\begin{Exercise}[#1, label=#4]
		#2
	\end{Exercise}
	\columnbreak
		#3
	\end{multicols}
	\begin{Answer}[ref=#4]
		#5
	\end{Answer}
}

% version 2
%\newcommand{\exemulticols}[5]{
%	\begin{Exercise}[#1, label=#4]
%	\begin{multicols}{2}
%		#2
%	\columnbreak
%		#3
%	\end{multicols}
%	\end{Exercise}
%	\begin{Answer}[ref=#4]
%		#5
%	\end{Answer}
%}

\newcommand{\exeTF}[1]{}

%%%%%%%%%%%%%%%%%%%%%%%%%%%%
% OPERATORS
%%%%%%%%%%%%%%%%%%%%%%%%%%%%

% deliminators
\DeclarePairedDelimiter{\abs}{\lvert}{\rvert}
\DeclarePairedDelimiter{\norm}{\lVert}{\rVert}
\DeclarePairedDelimiter{\scal}{\langle}{\rangle}

\DeclarePairedDelimiter{\ceil}{\lceil}{\rceil}
\DeclarePairedDelimiter{\floor}{\lfloor}{\rfloor}
\DeclarePairedDelimiter{\round}{\lfloor}{\rceil}

%  - others
\DeclareMathOperator{\Lap}{\mathcal{L}}
\DeclareMathOperator{\Var}{Var} % variance
\DeclareMathOperator{\Cov}{Cov} % covariance
\DeclareMathOperator{\proj}{proj} % projection
\DeclareMathOperator{\sym}{sym} % symétrique

\DeclareMathOperator{\RE}{Re} % partie réelle
\DeclareMathOperator{\IM}{Im} % partie imaginaire

\DeclareMathOperator{\pgcd}{pgcd}
\DeclareMathOperator{\ppcm}{ppcm}

\DeclareMathOperator{\ev}{ev}

\DeclareMathOperator{\tr}{tr}
\DeclareMathOperator{\sgn}{sgn}

\DeclareMathOperator{\Frac}{Frac}

% Lois
\DeclareMathOperator{\Bern}{\text{Bern}}
\DeclareMathOperator{\Binom}{\text{Binom}}

% I prefer the slanted \leq
\let\oldleq\leq % save them in case they're every wanted
\let\oldgeq\geq
\renewcommand{\leq}{\leqslant}
\renewcommand{\geq}{\geqslant}


%%%%%%%%%%%%%%%%%%%%%%%%%%%%
% commands
%%%%%%%%%%%%%%%%%%%%%%%%%%%%

% tel que
\newcommand{\tqs}{\text{ tels que }}
\newcommand{\tq}{\text{ tq. }}
\newcommand{\et}{\text{ et }}
\newcommand{\ou}{\text{ ou }}
\newcommand{\pourtout}{\text{ pour tout }}
\newcommand{\sct}{\text{ sachant }}
\newcommand{\id}{\text{id}}
\newcommand{\apriori}{\emph{a priori}}


% ensemble avec bigl et bigr
\newcommand{\bigset}[1]{\bigl\{ #1 \bigr\}}
\newcommand{\Bigset}[1]{\Bigl\{ #1 \Bigr\}}
\newcommand{\bigpar}[1]{\bigl( #1 \bigr)}
\newcommand{\Bigpar}[1]{\Bigl( #1 \Bigr)}

% PLUS INFTY AND MINUS INFTY WITH NO SPACE
\newcommand{\pinfty}{{+}\infty}
\newcommand{\minfty}{{-}\infty}

% vecteur flèche
\renewcommand{\vec}[1]{\overrightarrow{#1}}

% vecteur pmatrix
\newcommand{\pvec}[2]{\begin{pmatrix} #1 \\ #2 \end{pmatrix}}

% vecteur norme
%\newcommand{\norm}[1]{\left\Vert #1 \right\Vert}

% point plan
\newcommand{\point}[3]{
	#1\left(#2 ; #3 \right)
}

% \smash avant \underline pour coller la ligne au mot
% attention : ça empêche les line breaks donc c'est pas super...
\let\oldunderline\underline
\renewcommand{\underline}[1]{\oldunderline{\smash{#1}}}

% emph + index
% for \uline which allows underline with line breaks
\usepackage[normalem]{ulem}
% i add underline just for now
\newcommand{\emphindex}[1]{\emph{\uline{#1}}\index{#1}}

% tableau croisé
\newcommand{\tableaucroise}[4]{
\begin{tabular}{|c|c|c|c|}
	\cline{2-4}
	\multicolumn{1}{c|}{} & #1 \\ \hline
	#2 \\ \hline
	#3  \\ \hline
	#4  \\ \hline
\end{tabular}
}

% emdash in mathmode
\newcommand{\dash}{\text{---}}

%%%%%%%%%%%%%%%%%%%%%%%%%%%%
% ENVIRONMENTS
%%%%%%%%%%%%%%%%%%%%%%%%%%%%

% two columns toc
\newcommand{\twocolstableofcontents}{
	\renewcommand*\contentsname{}
	\begin{center}  
		{\Large \bfseries TABLE DES MATIÈRES \bigskip}    
	\end{center}
	\hrule
	\begin{multicols}{2}\setlength{\columnseprule}{1pt}
		\titleformat{\chapter}{}{}{0pt}{} 
		\titlespacing*{\chapter}{0pt}
		{*7} %  % vertical space before the title <<<<<<<<<<
		{*-15}  %  idem after title (in ex units + glue) <<<<<<<<<<<
		\tableofcontents
	\end{multicols}
	\hrule
}

% python minted
\newcommand{\python}[1]{
\inputminted[
		linenos,
		gobble=0,
		breaklines=true, % otherwise it breaks for no apparent reason?
		breakafter=,,
		fontsize=\normalsize,
		numbersep=8pt,
		tabsize=4, % tab ident = 4 spaces
		fontfamily=courier, %important pour les signes <, >
]{python}{python/#1.py}
}


% add points inside environments
% example : \begin{enumerate} \item \points{2} blabla \item \points{1} bla \end{enumerate}
\newcommand{\points}[1]{
	\leavevmode\marginpar{\makebox{[#1]}}\ignorespaces%
	\addtocounter{points}{#1}%
}

%%%%%%%%%%%%%%%%%%%%%%%%%%%%
% PACKAGE-SPECIFIC ENVIRONMENTS
%%%%%%%%%%%%%%%%%%%%%%%%%%%%

% global settings for the first level in enumerate environments
\setlist[enumerate, 1]{%
	align = left, labelwidth = \parindent%, labelsep = 0pt, leftmargin = \parindent
}


%%%%%%%%%%%%%%%%%%%%%%%%%%%%%%
% EXERCISES
%%%%%%%%%%%%%%%%%%%%%%%%%%%%%%


% counter for points (title of exercices)
\usepackage{totcount}
\newtotcounter{points}

\usepackage{ifthen}
\renewcommand{\ExerciseHeader}{
	\tikz[baseline=(R.base)]\node[draw,rectangle, thick, inner sep=2pt](R) {%
		\textbf{%
			\refAnswer{\ExerciseLabel}.%
		}%
	};\!
	%
	\ifnum\ExerciseDifficulty=0
	\else
		(\theExerciseDifficulty)
	\fi
}
\renewcommand{\DifficultyMarker}{$\star$}
\renewcommand{\AnswerHeader}{
	\tikz[baseline=(R.base)]\node[draw,rectangle, thick, inner sep=2pt](R) {\textbf{\theExercise.}};\!
}
\renewcommand{\exe}[4]{
	\begin{Exercise}[title=#1, label=#3]
		\if\relax\detokenize\expandafter{\ExerciseTitle}\relax
		%\marginpar{[Bonus]}
		\else
		\marginpar{\mbox{[\ExerciseTitle]}}
		\addtocounter{points}{\ExerciseTitle}
		\fi
		#2
	\end{Exercise}
	\begin{Answer}[ref=#3]
		#4
	\end{Answer}
}
\renewcommand{\exemulticols}[5]{
	\begin{multicols}{2}
	\begin{Exercise}[title=#1, label=#4]
		\if\relax\detokenize\expandafter{\ExerciseTitle}\relax
		%\marginnote{[Bonus]}
		\else
		\marginnote{\mbox{[\ExerciseTitle]\qquad}}
		\addtocounter{points}{\ExerciseTitle}
		\fi
		#2
	\end{Exercise}
	\columnbreak
		#3
	\end{multicols}
	\begin{Answer}[ref=#4]
		#5
	\end{Answer}
}





\SetDate[6/10/2026]
\reversemarginpar
\setlength{\marginparsep}{.5cm}

\begin{document}
\fancyhead[C]{\textbf{MG303 --- Contrôle Continu 1}}

%\null\vspace{-20pt}
%Nom : \hrulefill ~ou Code Apprenant :  \rule{3cm}{.1pt}

\begin{itemize}[label=$\bullet$]
	\item
	Matériel autorisé : aucun (pas de polycopié, pas de calculatrice).
	\item
	Les résultats des chapitres 1 et 2 peuvent être utilisés sans les redémontrer.
	Néanmoins, il n'est pas possible de faire référence aux exercices vus en classe ni au chapitre 3 (lemme de Hamza, par exemple) :
	si un problème ressemble à un exercice déjà vu, il faut le refaire.
	\item
	L'évaluation contient \ref{exe:last} problèmes.
	Les problèmes sont ordonnés par ordre croissant de difficulté. Le dernier problème est difficile.
	\item
	Chaque problème rapporte 5 points.
	Pour les questions vrai/faux, 1 point est accordé pour la bonne réponse, et 4 points pour la preuve ou le contre-exemple.
	\item
	L'ensemble $\K$ est un corps et, pour $p\in\N$ premier, l'ensemble $\F_p$ est le corps $\bigpar{\Z/p\Z, +, \cdot}$.
\end{itemize}

\hrule 

%\marginpar{[pts]}




\exe{}{
	Soit $p\in\K[X]$.
	Si $p$ n'admet aucune racine dans $\K$, alors $p$ est irréductible.


	Vrai ou faux ? Si vrai, donner une preuve. Sinon, donner un contre-exemple.
}{exe:VF1}{
	C'est faux.
	
	Le polynôme $(X^2 +1)^2$ dans l'anneau $\R[X]$ donc une contre-exemple suffisant.
	
	Attention en fait que ce n'est pas un contre-exemple dans $\C[X]$ !
}



\exe{}{
	Le polynôme $X^2 + 1$ est irréductible dans $\F_{17}[X]$.

	Vrai ou faux ? Si vrai, donner une preuve. Sinon, donner un contre-exemple.
}{exe:VF6}{
	C'est faux.
	
	En $x=4$, nous obtenons $4^2 + 1 = 0 \mod 17$.
	Il suit donc que $4$ est racine du polynôme et que $X-4$ divise $X^2 + 1$ dans $\F_{17}[X]$.
	En fait, par parité, $-4$ est aussi racine et 
		\[ X^2 - 1 = (X-4)(X+4). \]
}

\exe{}{
	Soit $p\in\K[X]$.
	Si $p$ est irréductible et $\deg(p)\geq2$, alors $p$ n'admet aucune racine dans $\K$.

	Vrai ou faux ? Si vrai, donner une preuve. Sinon, donner un contre-exemple.
}{exe:VF2}{
	C'est vrai et c'est une direction du lemme de Hamza.
	
	Si $p$ admet une racine $a$ dans $\K$, alors la division euclidienne de $p$ par $x-a$ montrer que $x-a$ divise $p$.
	Or, comme $\deg(p)\geq2$, cette division donne une réduction de $p$ : $p(x) = (x-a)q(x)$ avec $q$ non unité.
}





\exe{}{
	Soit $p\in\Z[X]$.
	Si $p$ est primitif, alors $p$ est irréductible dans $\Z[x]$.

	Vrai ou faux ? Si vrai, donner une preuve. Sinon, donner un contre-exemple.
}{exe:VF2.1}{
	C'est faux.
	
	Le produit de deux primitifs est primitif d'après le cours.
	Il suffit donc de prend $(X^2 + 1)^2 = X^4 + 2X^2 + 1$, par exemple.
	Ce polynôme est primitif mais se réduit dans $\Z[X]$.
}



\exe{}{
	L'ensemble $\Qsqrt2 = \bigset{ p + q\sqrt2 \tq p, q\in \Q } \subseteq \R$ est un corps lorsque muni de l'addition et de la multiplication usuelles des nombres réels.
	
	Vrai ou faux ? Si vrai, donner une preuve. Sinon, donner un contre-exemple.
}{exe:VF10-qs2}{
	C'est vrai.
	
	La seule partie difficile est de montrer que si $(p, q) \neq (0, 0)$, alors $p+\sqrt2q$ est inversible dans le corps.
	Or, comme $p^2 - 2q^2 \neq 0$ car $\sqrt2\not\in\Q$,
		\[ \frac1{p+\sqrt2q} = \frac{p-\sqrt2q}{p^2 - 2q^2} = \frac{p}{p^2 - 2q^2} - \frac{q}{p^2 - 2q^2} \sqrt2 \in\Qsqrt2. \]
}



\exe{}{
	Calculer le PGCD de $2X^3 - 5X^2 + 5X + 4$ et $-4X^3 +7X + 3$ dans $\Q[X]$.
}{exe:div-eucl-R}{
	Nous obtenons $X + \frac12$ après trois divisions euclidiennes.
}


\exe{}{
	Soient $a, b, c\in\K[X]$ trois polynômes.
	Montrer que si $a|(bc)$ et $\pgcd(a,b) = 1$, alors $a|c$.
}{exe:lem:Gauss}{
	Comme $\pgcd(a,b)=1$, il existe $u, v\in\K[x]$ tels que
		\[ au + bv = 1. \]
	En multipliant par $c$, nous obtenons que
		\[ c = a(uc) + (bc)v \]
	est divisible par $a$.
}


\exe{}{
	Soit $\alpha\in\C$ un nombre complexe dit « algébrique », c'est-à-dire qu'il existe un polynôme rationnel $f\in\Q[X]$ annulateur de $\alpha$ : $f(\alpha)=0$.
	
	Notons $m\in\Q[X]$ un « polynôme minimal » de $\alpha$, c'est-à-dire un polynôme non nul de $\Q[X]$ de \underline{degré minimal} vérifiant
	\begin{enumerate}[leftmargin=60pt]
		\item 
		$m$ est unitaire ; et
		\item
		$m$ est annulateur de $\alpha$ : $m(\alpha) = 0$.
	\end{enumerate}
	
	Montrer d'abord que $m$ divise dans $\Q[X]$ tous les polynômes rationnels $f$ annulateurs de $\alpha$.
	
	En déduire ensuite que le polynôme minimal est unique.
}{exe:VF10}{
	Soit $f\in\Q[X]$ un polynôme rationnel annulateur de $\alpha$.
	La division euclidienne de $f$ par $m$ donne, dans $\Q[X]$,
		\[ f = mq + r, \]
	avec $\deg(r) < \deg(m)$.
	Or, en $x=\alpha$, nous avons $r(\alpha)=0$.
	La minimalité du degré de $m$ implique donc que $r = 0$ et que $m$ divise $f$ dans $\Q[X]$.
	
	Nous concluons que $m$ est unique car deux polynômes minimaux $m$ et $m'$ unitaires et de même degré se divisent l'un l'autre.
	Nous avons donc $m= m' q$ et $q$ est constant égal à 1 pour respecter les degrés et le coefficient dominant égal à $1$.
}


%\exe{}{
%	Soit $\bigpar{A, +, \cdot}$ un anneau commutatif et $u, v\in A^\times$.
%	Alors $u\cdot v \in A^\times$.
%	
%	Vrai ou faux ? Si vrai, donner une preuve. Sinon, donner un contre-exemple.
%}{exe:VF8}{
%
%}




%\exe{}{
%	Le polynôme $3X^2 - 6$ est irréductible dans $\Z[X]$.
%
%	Vrai ou faux ? Si vrai, donner une preuve. Sinon, donner un contre-exemple.
%}{exe:VF5}{
%
%}


%\exe{}{
%	%La division euclidienne existe dans $\Z[X]$.
%	L'anneau $\Z[X]$ est un anneau principal.
%
%	Vrai ou faux ? Si vrai, donner une preuve. Sinon, donner un contre-exemple.
%}{exe:VF3}{
%
%}


%\exe{}{
%	Soient $a, b, c\in\K[x]$ trois polynômes.
%	Montrer que si $a|c$, $b|c$, et que $\pgcd(a, b)=1$, alors $(ab)|c$.
%}{exe:VF3}{
%	Notons $c = bq$ avec $q\in\K[x]$.
%	Comme $\pgcd(a,b)=1$, il existe $u, v\in\K[x]$ tels que
%		\[ au + bv = 1. \]
%	En multipliant par $q$, nous obtenons que
%		\[ aqu + cv = q \]
%	est divisible par $a$.
%	Par conséquent, $c = bq = ba\tilde{q}$ pour un $\tilde{q}\in\K[x]$ et $(ab)$ divise $c$.
%}

\exe{}{
	Soit $A$ un anneau commutatif principal. Un idéal $I$ de $A$ est dit « maximal » dès que les seuls idéaux $J$ de $A$ vérifiant $I \subseteq J \subseteq A$ sont $I$ et $A$.
	
	Montrer l'équivalence suivante, où $a\in A$ n'est pas une unité de $A$, et $(a)$ dénote l'idéal engendré par $a$.
	Comme vu en classe, $(a) = a\cdot A$.
	\begin{align*}
		\text{$(a)$ est maximal} && \iff && \text{$a$ est irréductible}
	\end{align*}
}{exe:VF9}{
	\underline{Direction $\implies$.}
	
	Soit $a = bc$ une réduction de $a$ en éléments de $A$.
	Il s'agit de montrer que soit $b$, soit $c$ est une unité de $A$.
	
	Considérons l'idéal $(b)$ engendré par $b$.
	Alors $(a) \subseteq (b)$ car $a\in(b)$.
	Par maximalité de $(a)$, il suit que $(b) = (a)$ ou $(b) = A$.
	
	Dans la première alternative, nous avons que $b\in(a)$ et donc que $b=aq = b(cq)$ et il suit que $cq = 1$ et donc que $c$ est une unité.
	
	Dans la seconde hypothèse, nous avons que $1\in(b)$ et donc que $bq = 1$ et il suit que $b$ est une unité.
	
	\underline{Direction $\impliedby$.}
	
	Soit $(a) \subseteq I \subseteq A$ un idéal de $A$. 
	Nous savons que $I = (b)$car $A$ est un anneau principal.
	
	Comme $a\in(b)$, il suit que $a = bq$ avec $q\in A$ et, par irréducibilité de $a$, soit $b$ est unité et $(b) = A$, soit $q$ est unité et $b=aq^{-1} \in (a)$ et donc $(a) = (b)$.
}







%\exe{}{
%	Soient $p, q\in\K[X]$ et $(p), (q)$ les idéaux engendrés par $p$ et $q$ respectivement.
%	Alors
%	\begin{align*}
%		p | q && \iff && (q) \subseteq (p)
%		% \text{ dans $\K[x]$ }
%	\end{align*}
%
%	Vrai ou faux ? Si vrai, donner une preuve. Sinon, donner un contre-exemple.
%}{exe:VF4}{
%
%}


%\exe{}{
%	Soit $A$ un anneau. Alors l'anneau $A[X]$ est factoriel.
%	
%	Vrai ou faux ? Si vrai, donner une preuve. Sinon, donner un contre-exemple.
%}{exe:VF7}{
%	Faux. $A = \sfrac{\Z}{8\Z}$ vu en classe où $X^2 - 1 = (X+1)(X-1) = (X+3)(X-3)$.
%}

%\exe{}{
%	Soit $A$ un anneau et $I, J$ deux idéaux de $A$. Alors $I \cap J$ est aussi un idéal de $A$.
%	
%	Vrai ou faux ? Si vrai, donner une preuve. Sinon, donner un contre-exemple.
%}{exe:VF9}{
%
%}





%\exe{3}{
%	Soit $p\in\K[X]$.
%	Le polynôme $p$ est dit « premier » dès qu'il est non nul et que, pour tous $a, b\in\K[x]$, 
%		\[ p|(ab) \implies p|a \ou p|b.\]
%	Montrer que tout polynôme premier est irréductible.
%}{exe:prime-implies-irr}{
%
%}




% bien mais utilise thm fondamental chapitre 3
%\exe{}{
%	Soient $a, b\in\R[x]$ deux polynômes à coefficients réels.
%	Montrer que les trois propositions suivantes sont équivalentes.
%	\begin{enumerate}[label=(\roman*)]
%		\item $\pgcd(a, b) \neq 1$ dans $\R[x]$
%		\item $\pgcd(a, b) \neq 1$ dans $\C[x]$
%		\item $a$ et $b$ admettent au moins une racine complexe en commun
%	\end{enumerate}
%}{exe:gcd-roots-R-C}{
%
%	\underline{Implication (i) $\implies$ (ii).}
%	
%	Si $d\in\R[x]$ non 1 divise $a$ et $b$ dans $\R[x]$, alors $d$ divise également $a$ et $b$ dans $\C[x]$.
%	
%	\underline{Implication (ii) $\implies$ (iii).}
%	
%	Soit $d=\pgcd(a,b) \in \C[x]$. Comme $d$ est non constant par hypothèse, $d$ admet une racine complexe $z\in\C$ d'après le théorème fondamental de l'algèbre.
%	Ainsi $d(z) = 0$ et donc $a(z) = 0$ et $b(z) = 0$ car $d$ divise $a$ et $b$.
%	
%	\underline{Implication (iii) $\implies$ (i).}
%	
%	Notons $z\in\C$ une racine complexe en commun à $a$ et $b$.
%	Par contradiction, si $\pgcd(a, b) = 1$ dans $\R[x]$, le théorème de Bézout impliquerait l'existence de polynômes $u, v \in\R[x]$ tels que $au + bv = 1$.
%	En plongeant cette égalité dans $\C[x]$, évaluer en $x=z$ aboutit à la contradiction $0=1$. \Lightning
%}


\exe{}{
	Soit $f\in\K[X]$ un polynôme non constant de degré $n\in\N^*$.
	\mbox{Montrer l'équivalence suivante.}
	\begin{center}
	\begin{tabular}{ccc}
		$\Bigpar{ \text{ $f$ est réductible } }$ & $\iff$ & $\Bigpar{\text{ il existe $a, b\in\K[X]$ non nuls avec $\deg(a) < n, \deg(b) < n$ tels que $f|(ab)$ }}$
	\end{tabular}
	\end{center}
}{exe:last}{
	\underline{Direction $\implies$.}
	
	Si $f = ab$ avec $a, b$ non constants, alors $f$ divise trivialement $ab$ et $\deg(a) + \deg(b) = n$ implique que $\deg(a), \deg(b) < n$.
	
	\underline{Direction $\impliedby$.}
	
	Considérons l'ensemble $E$ des couples $(a, b)$ de polynômes \underline{non nuls} vérifiant $\deg(a), \deg(b) < n$ et $f|(ab)$.
	Cet ensemble n'est pas vide par hypothèse.
	Choisissons un couple de $E$ pour lequel $\deg(a)$ est minimal et montrons que $a$ divise $f$.
	
	La division euclidienne de $f$ par $a$ dans $\K[x]$ s'écrit
		\[ f = aq + r, \]
	avec $\deg(r) < \deg(a)$.
	
	Comme $r = f -aq$, le polynôme $f$ divise $rb = rf - (ab)q$.
	Si $r\neq0$, l'appartenance $(r,b) \in E$ contredit le choix de $a$.
	Il suit donc que $r = 0$ et nous avons trouvé un diviseur propre de $f$.
	En effet, $a$ ne peut pas être constant car $\deg(b) < n$ et $\deg(ab) \geq n$.
}

%%%%%%%%%%%

\label{lastpage}
\newpage
\fancyhead[C]{\textbf{Solutions}}
\shipoutAnswer

\end{document}
